\chapter{Sistema eléctrico, instrumentos, luces y aviónica}
\label{cap:mm-electrico}

\opening{Un sistema de 14 voltios, un tacómetro electrónico del que depende la seguridad y una consola tan compacta que cada cable nuevo puede interferir con los mandos.}

\section{Sistema eléctrico}
\fuente{MM, capítulos 14, 33 y 35}

\begin{precaucion}
La instalación de dispositivos eléctricos puede afectar la exactitud y fiabilidad del tacómetro electrónico.
\end{precaucion}

El sistema incluye un \textbf{alternador de 14~V y 60~A} (70~A opcional), relé de batería, unidad de control del alternador (ACU) y \textbf{batería de 14~V}. La batería va en una caja de fibra de vidrio, normalmente en el tubo inferior izquierdo del bastidor; opcionalmente en la nariz bajo la consola superior, y en algunos R22 HP en el compartimiento de equipaje del asiento izquierdo. Los disyuntores, del tipo ``empujar para rearmar'' y rotulados con función y amperaje, están en el reborde delante del asiento delantero izquierdo.

El interruptor maestro controla el relé de batería, que la desconecta de todos los circuitos salvo el tacómetro y el reloj, alimentados directamente desde la batería a través de un terminal del interruptor del embrague.

\textbf{Unidad de control del alternador}: detecta la tensión en el shunt del amperímetro mediante un cable de sensado remoto y cumple tres funciones:
\begin{itemize}
\item regula la salida para mantener la batería a \textbf{13,4--13,9~V};
\item enciende la luz ALT si la tensión cae a \textbf{12,55--12,95~V};
\item corta el campo del alternador por sobretensión a \textbf{15,75--16,25~V}.
\end{itemize}

En instalaciones sin A942-1, la tensión de salida (13,2--13,8~V en el terminal ``I'' del regulador con el motor en marcha) se ajusta con el tornillo de rango del regulador.

El circuito del actuador del embrague tiene, además de su disyuntor, un fusible de bajo amperaje (1,5~A) para que una sobrecarga del motor no dispare el disyuntor y apague prematuramente la luz del embrague. Las luces de aterrizaje están cableadas a través del interruptor del embrague y se apagan al desacoplarlo. La luz de mapa está en el circuito de instrumentos como respaldo ante un cortocircuito del circuito de luces.

\subsection{Batería}

El mantenimiento de la batería sigue el manual del fabricante (Concorde). \textbf{Al desconectar, primero el cable negativo y luego el positivo; al conectar, primero el positivo y luego el negativo.} Pernos de terminales con el torque indicado en la etiqueta de la batería.

\subsection{Actuador del embrague}

Tras el arranque, el motor se acopla a la transmisión mediante las correas en V, tensadas al elevar la polea superior con un actuador eléctrico situado entre las poleas. El actuador detecta la carga de compresión (tensión de correas) y se detiene al alcanzar la tensión correcta. La luz CLUTCH se enciende siempre que el circuito está energizado (acoplando, desacoplando o retensando) y permanece encendida hasta que las correas están bien tensadas o totalmente sueltas.

La holgura de las correas en el arranque debe permitir que \textbf{las palas empiecen a girar dentro de los 5 segundos} de conectado el embrague; una holgura excesiva puede hacer saltar las correas de las ranuras. Si el fusible se funde, el motor se detiene pero la luz queda encendida; con el disyuntor abierto no hay tensado ni luz.

\begin{precaucion}
Nunca despegue con la luz del embrague encendida.
\end{precaucion}

\subsection{Luces de advertencia y precaución}

Embrague, sobretemperatura de la caja principal, virutas en caja principal y de cola, arranque conectado, bajo combustible, bajas RPM, alternador, baja presión de aceite, freno rotor, gobernador desconectado, monóxido de carbono (con calefacción) y acelerador a fondo. La luz y la bocina de bajas RPM indican \textbf{97\,\% o menos}. Las de baja presión de aceite y bajo combustible tienen sensores propios, independientes de los indicadores. Los detectores de virutas son magnéticos, en el tapón de drenaje de cada caja; la sobretemperatura la acciona un interruptor térmico cerca del piñón de entrada.

\textbf{Luz de acelerador a fondo}: con la válvula de combustible cerrada y batería conectada, colectivo arriba, abra el acelerador hasta que se encienda: la separación entre el tope del acelerador y la esquina del tope de plena apertura del cuerpo debe ser \textbf{0,114--0,148~in}. Se ajusta girando la leva ranurada A607-2 (bajo el asiento del copiloto) para que el interruptor V3-1 actúe a unos 0,130~in.

\textbf{Detector de monóxido de carbono}: sensor sobre la salida de calefacción del piloto. Al conectar la energía la luz parpadea dos veces (autoprueba); un destello cada 4 segundos indica falla del sensor. Si la luz se enciende: cierre la calefacción, abra ventilaciones, salga del estacionario; con síntomas (dolor de cabeza, somnolencia, mareo), aterrice de inmediato e inspeccione el escape antes del siguiente vuelo. Evite solventes, detergentes y aerosoles cerca del sensor y tape sus aberturas al limpiar la cabina.

\textbf{Bocina de bajas RPM}: con el colectivo levantado \textbf{0,2--0,4~in} (en la empuñadura) sobre el tope inferior, la unidad debe activar bocina y luz a \textbf{97--96\,\%} y apagarlas por encima de 96--97\,\%. La unidad está dentro de la consola superior; se ajusta por un agujero de 3/8~in en el panel vertical izquierdo: A569-1 (circuito encapsulado) $\approx$ 1/16 de vuelta por 1\,\%; A569-5 (caja metálica) $\approx$ 2 vueltas por 1\,\%. Si no se logra, se reemplaza. Las aeronaves recientes tienen además una alerta de \textbf{altas RPM} (tono alternado al acercarse a 110\,\%) con botón de prueba.

\subsection{ELT}

Transmisor con batería interna en el bastidor superior (acceso por la puerta superior trasera del capó), antena externa e interruptor remoto con indicador a la izquierda del cíclico. Ambos interruptores en \textbf{AUTO/ARM} para el vuelo; ON para activarlo ante un aterrizaje de emergencia inminente; RESET para detener una activación accidental. Las pruebas siguen las instrucciones del fabricante.

\subsection{Unidad de monitoreo del motor (EMU)}

En aeronaves recientes, dentro de la caja del gobernador D270-1, registra RPM de motor y rotor, temperatura de aceite y CHT. Luz ámbar en el compartimiento de equipaje derecho: tras unos 10~s de autoprueba, \textbf{fija = normal}; \textbf{parpadeo lento} (cada 2~s) o apagada = falla de sensores o circuitos; \textbf{parpadeo rápido} (4 por segundo) = excedencia. Una falla o excedencia debe ser investigada y reseteada por un mecánico calificado antes del siguiente vuelo (la app EMU para iPad permite verla, no resetearla). Es solo una ayuda de mantenimiento: reportar excedencias sigue siendo responsabilidad del piloto.

\subsection{Otros sistemas}

\begin{itemize}
\item \textbf{Protección de sobrevelocidad en el arranque} (S/N 4825 y posteriores): con el embrague desacoplado, si el motor supera 85--90\,\% con el rotor por debajo del 50\,\%, el tacómetro o el gobernador activan un relé que pone a masa los cables P de ambos magnetos. No requiere mantenimiento periódico ni evita todas las sobrevelocidades.
\item \textbf{Encendido electrónico Lycoming (EIS)}: en aeronaves recientes reemplaza al magneto izquierdo de arranque y al vibrador; véase Lycoming SI 1569.
\item \textbf{Cámara de cabina} opcional: graba video 4K, audio y GPS en una memoria extraíble; no retire la memoria mientras graba.
\item \textbf{Sensores}: termistor de aceite A760 a 300~in-lb (terminal 20~in-lb) frenado con alambre de 0,032~in; sensor CHT en el cilindro \#4 a 70--80~in-lb; sensor OAT A760-4 a 18~in-lb; sensor de efecto Hall C143-1 a 96~in-lb.
\item \textbf{Calibración del CHT}: con una resistencia de 32 $\pm$ 0,4~$\Omega$ (equivale a 500~°F) en lugar de la sonda, ajuste el potenciómetro de la placa C691-1 hasta que la aguja quede entre ``tocando el rojo'' y ``media aguja en el rojo'', y fíjelo con cinta de aluminio.
\end{itemize}

\subsection{Diagnóstico eléctrico}

\begin{center}
\small
\begin{tabularx}{\linewidth}{@{}>{\raggedright\arraybackslash}p{0.24\linewidth}>{\raggedright\arraybackslash}p{0.38\linewidth}>{\raggedright\arraybackslash}X@{}}
\toprule
Falla & Causa probable & Corrección \\
\midrule
Sin energía & Bornes corroídos, mala masa, disyuntor disparado, batería baja, alternador sin salida, cable o terminal defectuoso & Limpiar, verificar, recargar, revisar correa y ACU, reemplazar \\
Arranca lento y no prende & Batería baja; poca holgura de correas; bornes sucios; relé de arranque & Cargar; ajustar tornillo de límite inferior del actuador; limpiar; reemplazar \\
Gira pero no arranca & Llave de encendido, vibrador de arranque, retardo del magneto & Reemplazar o ajustar \\
El arranque no funciona & Freno rotor aplicado; batería baja; disyuntores HORN START y CLUTCH START disparados; actuador no desacoplado del todo; conexiones & Soltar freno; cargar; rearmar; acoplar y desacoplar del todo; revisar cableado \\
Retroceso del arrancador & Contacto de retardo del magneto izquierdo desplazado (debe medir 0,609 $\pm$ 0,10~in con leve rebote); cableado del vibrador; llave (SL-29); acoplamiento de impulso & Corregir; kit KI-4 \\
Batería descargada o de vida corta & Batería agotada, régimen de carga incorrecto, equipo encendido, impurezas & Reemplazar, ajustar regulador, recargar \\
Alternador no carga & Alternador, regulador, relé de sobretensión o interruptor & Reemplazar \\
\bottomrule
\end{tabularx}
\end{center}

\textbf{Actuador del embrague}: si no acopla o no desacopla, revise motor agarrotado, fusible de 1,5~A, circuito abierto, interruptor de sobrecarrera disparado, tensión en el disyuntor, interruptores de resorte o de límite superior y masa del cable -88. Luz que parpadea en vuelo: resorte de columna descalibrado, correas A190 no emparejadas, polea inferior mal apretada o rodamiento A184/A181 áspero. \textbf{Luz encendida 1--6~s en vuelo: normal} (retensado). \textbf{Más de 6~s}: correas estiradas más allá del interruptor de sobrecarrera (reemplazarlas) o interruptor accionado por una fuerza externa.

\subsection{Valores de calibración}

\begin{center}
\small
\begin{tabular}{lcc}
\toprule
Tacómetro A792 & RPM motor & RPM rotor \\
\midrule
97\,\% & 2471--2476 & 494--495 \\
100\,\% & 2548--2552 & 509--510 \\
104\,\% & 2649--2655 & 529--530 \\
\bottomrule
\end{tabular}
\end{center}

Se verifican con equipo Vibrex 2000 o similar (resolución $\pm$1~rpm, calibrado en el último año). Solo puede ajustarse el lado del rotor: con el motor a 104\,\%, gire el tornillo trasero (1/8 de vuelta $\approx$ 1\,\%; horario aumenta) hasta que el rotor indique 104\,\%. Las agujas deben coincidir dentro del 1\,\%. El A792-1 antiguo no es ajustable.

\begin{center}
\small
\begin{tabularx}{\linewidth}{@{}>{\raggedright\arraybackslash}p{0.3\linewidth}>{\raggedright\arraybackslash}X@{}}
\toprule
Instrumento & Calibración \\
\midrule
Combustible (6246-00473) & 0,70~$\Omega$ = E; 21,2~$\Omega$ = 1/4; 42,0~$\Omega$ = 1/2; 67,5~$\Omega$ = 3/4; 90,0~$\Omega$ = F ($\pm$1 ancho de aguja). Potenciómetros ``Null'' (vacío a 0,7~$\Omega$) y ``Gain'' (mitad a 42~$\Omega$) bajo cinta de aluminio. \\
Amperímetro (6246-00468) & $\pm$104~mV = $\pm$70~A; $\pm$52~mV = $\pm$35~A; 0~mV = 0~A \\
CHT (6246-00088) & 745~$\Omega$ = 200~°F; 110~$\Omega$ = 350~°F; 34~$\Omega$ = 500~°F \\
Presión de aceite (6246-00647, línea roja 115~psi) & 9~$\Omega$ = 0; 46~$\Omega$ = 25; 84~$\Omega$ = 55; 131~$\Omega$ = 95; 152~$\Omega$ = 115~psi \\
Temperatura de aceite (6246-00090) & 903,5~$\Omega$ = 75~°F; 100~$\Omega$ = 180~°F; 36~$\Omega$ = 245~°F \\
Temperatura del carburador (A604-2) & 77,40~$\Omega$ = $-$30~°C; 90,38~$\Omega$ = 0~°C; 103,36~$\Omega$ = 30~°C. Una sonda en hielo picado y agua debe medir 90,38~$\Omega$. Sonda a 3--4~in-lb (más la daña). \\
\bottomrule
\end{tabularx}
\end{center}

Los instrumentos del grupo se calibran inclinados 20--25° respecto de la horizontal.

\subsection{Carga eléctrica}

La carga continua máxima del alternador de 60~A es \textbf{54~A}. La carga de una aeronave se calcula sumando las cargas continuas e intermitentes del equipo instalado (tabla del MM, \S\,33-120), o se mide en el borne de la batería con el alternador apagado y todo el equipo encendido, escalando por la relación entre tensión de batería y nominal. Algunos consumos típicos: luces de aterrizaje incandescentes 2 $\times$ 7,7~A; HID 2 $\times$ 2,9~A (9,52~A al encender); luces de posición incandescentes 3 $\times$ 2,2~A; Aspen 4,8~A; GTN~650 2,65~A; soplador de calefacción 7,5~A; motor de arranque 200~A intermitente.

\begin{advertencia}
La instalación de equipo electrónico fuera de fábrica puede ser peligrosa. La consola y el túnel de mandos son tan compactos que cualquier cable adicional puede interferir con los controles de vuelo, y los tacómetros, sistemas de advertencia y equipos de navegación son sensibles a interferencias. Un dispositivo no probado ni aprobado por RHC puede crear una condición peligrosa.
\end{advertencia}

Los diagramas de cableado completos están en los capítulos 14 y 35 del MM y no se reproducen aquí.

\section{Instrumentos}
\fuente{MM, capítulo 13}

Instrumentos estándar: anemómetro, tacómetro doble de motor y rotor, altímetro, manifold, brújula; amperímetro, presión y temperatura de aceite, CHT, cantidad de combustible; reloj, temperatura del carburador y OAT digital. El horómetro, accionado por la presión de aceite, está delante del asiento del piloto. Con pantallas electrónicas puede usarse el filtro D327 para reducir reflejos nocturnos.

\begin{precaucion}
Proteja los instrumentos con espuma: trátelos ``como huevos''. No use herramientas magnetizadas cerca del tacómetro ni de las antenas.
\end{precaucion}

\subsection{Sistema pitot-estático}

El tubo pitot está en el borde delantero del carenado del mástil; la toma estática, dentro del capó trasero junto a la bisagra de la puerta. El agua se drena retirando los tapones plásticos accesibles por el panel de inspección inferior de la cabina (solo si anemómetro o altímetro son erráticos). Inspeccione con frecuencia que no haya insectos u obstrucciones.

\begin{precaucion}
Nunca aplique succión al pitot ni presión a la estática; aplicar o soltar la succión bruscamente puede dañar anemómetro, altímetro y variómetro.
\end{precaucion}

\begin{center}
\small
\begin{tabularx}{\linewidth}{@{}>{\raggedright\arraybackslash}p{0.22\linewidth}>{\raggedright\arraybackslash}p{0.44\linewidth}>{\raggedright\arraybackslash}X@{}}
\toprule
Prueba & Método & Fuga máxima \\
\midrule
Pitot & Tape el drenaje del codo; con una jeringa, presión lenta hasta 70~kt; golpee suavemente el vidrio & 10~kt en 1~min \\
Estática & Altímetro a la elevación local; succión lenta hasta +500~ft & 100~ft en 1~min \\
Manifold & Desconecte la línea A740-1 del motor; succión hasta 10~inHg & sube 1~inHg en 1~min \\
\bottomrule
\end{tabularx}
\end{center}

\subsection{Verificaciones de instrumentos}

\begin{itemize}
\item \textbf{Variómetro}: cero en estacionario.
\item \textbf{Anemómetro}: solo correcto en vuelo hacia adelante; hasta 102~kt (línea roja).
\item \textbf{Altímetro}: con el reglaje local, debe coincidir con la elevación del campo dentro de \textbf{70~ft}.
\item \textbf{Manifold}: con el motor parado, dentro de \textbf{0,3~inHg} de la presión barométrica.
\item \textbf{Temperatura del carburador}: con motor frío, aproximadamente igual a la OAT.
\item \textbf{CHT}: termopar en el cilindro \#4.
\item \textbf{Sonda OAT}: no se permite cortar ni empalmar sus cables (longitud calibrada).
\item \textbf{Brújula}: compensación en rosa de vientos, a RPM gobernadas, cada 30°, con toda la aviónica, luces exteriores y dimmer al máximo: \textbf{error máximo 5°}; registre en la tarjeta y verifique los 4 cardinales con luces apagadas y solo COM1 y transpondedor.
\item \textbf{Amperímetro inoperativo}: verifique los dos fusibles de 3~A en el conjunto B304-18 del shunt. \textbf{Reloj}: fusible AGC-3 en el cortafuego.
\item \textbf{Aspen PFD}: ayuda de conciencia situacional junto con los instrumentos VFR requeridos; disyuntor EFIS 7,5~A; la \textbf{batería interna se prueba cada 12 meses}; tras reinstalarlo, prueba de fugas pitot-estática.
\end{itemize}

Ante indicaciones nulas o erráticas, el MM sigue el mismo patrón: buscar obstrucciones, agua, líneas sueltas, dobladas o rotas, hacer las pruebas de fugas y, si persiste, devolver el instrumento a RHC.

\section{Luces}
\fuente{MM, capítulo 24}

Luz anticolisión roja en el cono de cola (interruptor STROBE, disyuntor 5~A); luces de posición a cada lado de la cabina y en la cola (NAV LTS, disyuntor LTS 5~A); iluminación de instrumentos por luces de poste (antiguas) o una luz en la parte superior del parabrisas (recientes), regulada con la perilla sobre el interruptor de luces; luz de mapa giratoria en el techo, utilizable como iluminación de emergencia del panel. Dos \textbf{luces de aterrizaje} en la nariz con distinto ángulo, con un solo interruptor; \textbf{solo funcionan con el interruptor del embrague en ``acoplado''} (disyuntor LAND LT 15~A con HID, 20~A con lámparas B196-26). Se recomienda el uso continuo de luces de aterrizaje y posición en vuelo para evitar colisiones.

Al reinstalar la anticolisión rojo/blanca, el lente rojo va hacia adelante.

\section{Aviónica}
\fuente{MM, capítulo 34}

\begin{itemize}
\item Desde S/N 4591, 4650, 4651, 4656 y posteriores hay \textbf{interruptor de aviónica} (relé ``a prueba de fallas'': falla en ON); antes, un relé maestro de radio cortaba la barra de aviónica durante el arranque.
\item \textbf{Transpondedor Garmin GTX 330 ES} (ADS-B Out por \emph{squitter} extendido de 1090~MHz, con posición del GPS principal; requerido en ciertos espacios aéreos de EE.\,UU.\ desde el 1 de enero de 2020). El sistema ADS-B del R22 requiere un GPS Garmin GTN serie 600 (solo en la consola del lado del piloto). Un cambio de matrícula puede requerir reprogramación.
\item \textbf{COM Garmin GTR 225B} (14/28~V) y \textbf{audio Garmin GMA 350H} (audio 3D, comandos de voz).
\item Los equipos Garmin instalados en fábrica son \textbf{``on condition''}: sin mantenimiento periódico; tienen autoprueba al encender. La inspección especial consiste en revisar equipo, bandeja, barras de cobre, disyuntor y cableado (sin cables sueltos, rozados o rotos ni señales de arco).
\item \textbf{Antenas}: quite pintura e imprimación en las superficies de contacto para asegurar la masa; cordón de sellador B270-1 de máx.\ 0,1~in de alto en las esquinas.
\item \textbf{Software}: solo versiones y configuraciones aprobadas por la FAA para Robinson; no actualice sin conocer su estado de aprobación y compatibilidad (puede requerir actualizar otros equipos). Mientras el equipo funcione bien, no hay requisito de verificar o actualizar; las actualizaciones obligatorias llegan por SB o AD. Las bases de datos (cartas, terreno) sí pueden actualizarse regularmente.
\end{itemize}
